\section{Related Work}
\label{sec:related}

\paragraph{Incentives and pricing in data markets.}
Designing incentives for data exchange is a long-standing concern~\citep{liu2024data,azcoitia2022survey}. One line of work sells data or statistics through auctions, often under privacy constraints~\citep{ghosh2011selling,nissim2014redrawing,zhang2020selling,agarwal2024towards}, while another studies pricing of data and machine-learning products~\citep{chen2017model,chen2019towards,mi2025multi,qin2023research}. Incentive mechanisms have also been developed for federated learning and collaborative analytics~\citep{yu2020fairness,li2023martfl,han2024dynamic,zhang2025incentive}, and for revenue or profit sharing among contributors~\citep{cao2017game,bauer2024designing}. These works largely treat data as a good that is consumed by its buyer, and do not model the strategic consequences of buyers becoming sellers themselves.

\paragraph{Replicability and resale of data.}
Unlike rival goods, data is non-rival and can be replicated and resold~\citep{JonesTonetti2020NonrivalryData,shiller2013digital,ali2024reselling}. A handful of works acknowledge that resale can generate value along a transaction chain~\citep{biswas2021incentive,xin2025tbds,ma2026incentivizing}.
\citet{biswas2021incentive} design a marketplace that permits resale but do not analyze incentive design for it.
\citet{xin2025tbds} propose a blockchain-based mechanism, \textsc{TBDS}, that traces transactions from resale and pays upstream contributors a fixed fraction of downstream revenue.
\citet{ma2026incentivizing} first studies the benefits of incentive designs on data trading with data resale. However, their analysis is only preliminary and relies on strict assumptions that data owner can sell data to at most one buyer.


\paragraph{Mechanisms on networks and sequential trade.}
Our reallocation of value to upstream participants is reminiscent of diffusion auctions, referral mechanisms, and multi-level marketing, which reward intermediaries for propagating an item or information through a network~\citep{emek2011mechanisms,li2017mechanism,li2022diffusion,zhang2020redistribution}. A separate line studies equilibrium behavior in sequential bilateral or intermediated trade~\citep{condorelli2017bilateral,manea2018intermediation,sandholm2006sequences}. Our setting differs in three respects: the traded object is replicable data, so a seller does not relinquish it; the focus is on how an external reallocation mechanism reshapes equilibrium trading behavior; and trades occur on a platform among mutually anonymous players via posted prices, rather than on a known social network.

% \paragraph{Positioning of this work.}
% The closest works are \citet{xin2025tbds} and \citet{ma2026incentivizing}. We depart from both in two ways. First, we move from a chain to a \emph{tree}, modeling the pervasive one-to-many resale in which a data product is sold to multiple buyers; the chain is special case of trees when each node has single child. Second, whereas \citet{xin2025tbds} fix the reallocation ratio and \citet{ma2026incentivizing} analyze the chain, we treat \emph{any} budget-feasible reallocation rule on the tree, give polynomial-time and FPTAS algorithms to compute the induced equilibria, and prove that reallocation weakly expands trade and welfare in this richer topology. The tree introduces genuinely new difficulty---a seller balances many downstream subtrees and a local change in the mechanism propagates across branches---which our subtree decomposition and edge-by-edge comparison are designed to handle.
